\documentclass[10pt,twocolumn,letterpaper]{article}
\usepackage[margin=0.70in,top=0.70in,bottom=0.74in,columnsep=0.24in]{geometry}
\usepackage{amsmath,amssymb,bm,booktabs,array,tabularx}
\usepackage[colorlinks=true,linkcolor=black,citecolor=black,urlcolor=black]{hyperref}

\setlength{\parindent}{0pt}
\setlength{\parskip}{2.8pt}
\renewcommand{\arraystretch}{1.05}

\begin{document}

\twocolumn[{%
\begin{center}
{\large\bfseries A US Long-Short Equity Portfolio: Tail-Weighted Forecasts, Adverse-Tail Risk, and an Audit of What Did Not Work}\\[6pt]
{\normalsize QUARCC}\\
{\small Quantitative Research Note}\\[3pt]
{\small November 2026}
\end{center}
\vspace{2pt}
\begin{center}
\begin{minipage}{0.91\textwidth}
\small
\textbf{Abstract.} We document a monthly, dollar-neutral US equity portfolio and the research that followed it. Annual walk-forward Ridge and tail-weighted XGBoost models score 147 characteristics from Jensen, Kelly, and Pedersen (2023) across 529,082 stock-months. A convex program with an alpha-alignment factor, an adverse-tail penalty, and explicit trading and borrow costs forms the book. Replayed in an event-driven engine from \$10 million over January 2021 to August 2026 (68 months), the portfolio has a net information ratio of 3.77 against T-bill $+4\%$, 12.4\% volatility, a $-4.5\%$ maximum drawdown, and $18.19\times$ growth. We do not read these figures as an achievable return. Models are walk-forward, but the design was chosen after 65 trials that had seen the evaluation window (deflated Sharpe ratio 0.886). About 91\% of short gross sits in micro- and nano-capitalisation stocks, borrow availability is modelled and not observed, and the book's liquidity problems grow quickly with assets. Forecast skill sits almost entirely on the short side and appears in the median stock, not the mean. A correction to how delisted stocks are handled raises the headline ratio, but only because of an assumption we highlight. Later attempts to replace the design failed, and a registered synthetic test finds no forecast contribution in the most ambitious successor.
\end{minipage}
\end{center}
\vspace{6pt}
}]

\section{Data and Walk-Forward Protocol}

The panel holds 529,082 stock-months of US common equities from January 2015 to August 2026, with 147 characteristics from Jensen, Kelly, and Pedersen (2023; hereafter JKP). Each month, with information dated at or before the previous month-end, models score the universe. They are refit every January on all earlier months. The evaluation window is the 68 months from January 2021 to August 2026. Each monthly book is replayed in NautilusTrader from \$10 million with whole shares, Interactive Brokers fees, spread, impact, borrow, and recall.

\section{Forecast and Alpha}

\textbf{Blend.} Following Gu, Kelly, and Xiu (2020), the signal averages cross-sectional $z$-scores of a Ridge regression ($\alpha = 1000$) and a tail-weighted XGBoost model (400 trees, depth 6, Platt-calibrated, trained on names priced at \$1 or more):
\begin{equation}
z_{i,t} = \mathrm{clip}\Big(z_{\text{CS}}\big(\tfrac{1}{2} z_{\text{CS}}(\hat{y}_{i,t}^{\text{ridge}}) + \tfrac{1}{2} z_{\text{CS}}(\hat{y}_{i,t}^{\text{xgb}})\big), -3, 3\Big). \label{eq:blend_z}
\end{equation}
Candidates are the 200 highest and 200 lowest scores among names priced at \$1 or more, for which a 1\% position is at most 10\% of average daily volume and 5\% of market capitalisation.

\textbf{Alpha.} A conformalised quantile model (Romano, Patterson, and Cand\`es, 2019) sets a confidence weight, and alpha follows the fundamental law (Grinold and Kahn, 2000):
\begin{equation}
\alpha_{i,t} = \mathrm{IC}^+_{\text{leg}(i),t}\,\sigma_{i,t}\,z_{i,t}\,\mathrm{conf}_{i,t}, \label{eq:gk_alpha}
\end{equation}
where $\mathrm{IC}^+_{\text{leg},t}$ is the leg's rank IC over the previous 36 months, floored at zero, and $\sigma_{i,t}$ is model volatility.

\textbf{The asymmetry that defines the book.} Within the 200-name short pool, rank IC is $+0.062$ in 2016--2020 and $+0.142$ in 2021--2026. Within the long pool it is $-0.012$ and $-0.011$. As a result $\mathrm{IC}^+_{\text{long},t}$ is zero in 82\% and 91\% of months respectively, and longs mainly hedge the factor, beta, and industry exposure of the shorts. In the evaluation window, shorts contribute $+44.8\%$ a year and longs $+9.8\%$. With a negative long-pool IC, long P\&L should not be read as evidence of selection skill.

\textbf{Median, not mean.} The bottom-quintile blend score predicts the median stock, about 1.3--1.4\% a month below the universe ($t = -3.5$), but the mean difference is zero ($t = 0.0$). Forecasts identify typical losers, and rare large winners in the same names offset them. Earning the gap means surviving that right tail, which is why position caps and tail penalties matter here. Platt calibration also fit an inverted slope in 2020, which is a warning about the stability of the XGBoost score.

\section{Risk and Portfolio Construction}

\textbf{Covariance.} A fundamental factor covariance $\bm{\Sigma}_{F,t}$ is augmented with an alpha-alignment factor (Ceria, Saxena, and Stubbs, 2012) along unit-norm alpha $\hat{\bm{\alpha}}_t$, rescaled by a variance correction $g_t \in [0.5, 4]$ learned from the book's realised and predicted variance, and given a diagonal adverse-tail term:
\begin{equation}
\bm{\Sigma}_t = g_t\big(\bm{\Sigma}_{F,t} + \sigma_{\text{AAF},t}^2 \hat{\bm{\alpha}}_t \hat{\bm{\alpha}}_t^\top\big) + \mathrm{diag}(\mathbf{h}_{\text{tail},t}). \label{eq:sigma}
\end{equation}
Here $h_{\text{tail},i,t}$ is built from the 0.99 conformal quantile of losses on the held side, namely squeezes for shorts and crashes for longs, in excess of the Gaussian quantile. It is scaled by a power-law conditional tail-variance multiplier (Taleb, 2020).

\textbf{Costs.} Half-spreads scale the high--low range (Corwin and Schultz, 2012) by a factor of 0.3947, calibrated on the most liquid names in 2016--2020. Impact follows Almgren et al.\ (2005), with terms in $|\Delta w_i|^{1.6}$ and $\Delta w_i^2$ evaluated at current NAV. Shorts pay a borrow fee tiered by dollar volume (D'Avolio, 2002) plus expected recall cost.

\textbf{Program.} Each month the book solves
\begin{align}
\max_{\mathbf{w}} \;\; & \bm{\alpha}_t^\top \mathbf{w} - \lambda m_t \mathbf{w}^\top \bm{\Sigma}_t \mathbf{w} - \mathrm{TC}_t(\Delta\mathbf{w}) \nonumber\\
& - \mathrm{Carry}_t(\mathbf{w}^-) - \tau \|\mathbf{w} - \mathbf{w}_{\text{drift}}\|_1, \label{eq:qp}
\end{align}
with $\lambda = 2$, drawdown multiplier $m_t = 1 + 40\,\mathrm{DD}_{t,6\text{m}}$, and $\tau = 0.01$. Long and short legs are each exactly 50\% of NAV. Beta on Fama--MacBeth up and down market betas is within $\pm 0.05$, the book may not be short market convexity (Ang, Chen, and Xing, 2006), Fama--French 49-industry net exposure is within $\pm 5\%$, and single names are capped at 2.5\%. We solve with Clarabel (Goulart and Chen, 2024) through CVXPY (Diamond and Boyd, 2016).

\section{Results}


Table~\ref{tab:results} reports the portfolio as an engine run. Costs are about 5.2\% a year: 1.8\% trading (commission, spread, and impact) and 3.4\% borrow and recall. Average holdings are 410 names (263 long, 146 short), with 4 of 68 months above 500 names because of small residual positions (maximum 509). Top-ten winners account for 22\% of gross P\&L across 1,847 names traded.

\begin{table*}[t]
\centering
\caption{Results, January 2021 to August 2026 (68 months), net of modelled costs, from the \$10 million ledger. Sharpe and alpha use monthly excess returns; alpha's $t$-statistic uses Newey--West errors. S\&P 500 figures use the price index.}
\label{tab:results}
\footnotesize
\setlength{\tabcolsep}{5pt}
\begin{tabular}{lrr}
\toprule
Metric & Portfolio & S\&P 500 \\
\midrule
Growth of \$1 & 18.19 & 2.05 \\
CAGR / arithmetic mean & 66.8\% / 53.0\% & 13.5\% / 13.8\% \\
Volatility & 12.36\% & 15.10\% \\
Sharpe ratio & 4.10 & 0.70 \\
Maximum drawdown & $-4.49\%$ & $-24.8\%$ \\
IR vs.\ T-bill $+4\%$ & 3.77 & 0.44 \\
Beta / alpha to S\&P 500 ($t$) & $+0.09$ / 48.9\% (8.35) & --- \\
Rolling 12-month beta & $-0.31$ to $+0.99$ & --- \\
Costs per year & 5.2\% & --- \\
Monthly one-way turnover & 13.9\% & --- \\
\bottomrule
\end{tabular}
\qquad
\begin{tabular}{lrr}
\toprule
Year & Portfolio & S\&P 500 \\
\midrule
2021 & 17.3\% & 26.9\% \\
2022 & 65.2\% & $-19.4\%$ \\
2023 & 76.4\% & 24.2\% \\
2024 & 97.7\% & 23.3\% \\
2025 & 106.5\% & 16.4\% \\
2026 (to Aug.) & 30.3\% & 12.3\% \\
\bottomrule
\end{tabular}
\end{table*}

\begin{table*}[t]
\centering
\caption{Attribution of the universe and delisting corrections. Engine runs from \$10 million, January 2021 to August 2026. All rows share parameters and rebalance rules. Maximum drawdown is on excess returns. The bold row is the portfolio reported above. The universe labelled ``all rows'' scores every stock at formation, including those whose next-month return is missing.}
\label{tab:audit}
\footnotesize
\setlength{\tabcolsep}{5pt}
\begin{tabular}{llrrrrr}
\toprule
Universe, exit liquidity & Missing return treatment & IR & Growth & Max DD & Beta & Costs/yr \\
\midrule
Original, original exits & $-30\%$ long, 0 short & 3.631 & 17.81 & $-5.1\%$ & $+0.118$ & 5.13\% \\
Original, repaired exits & $-30\%$ long, 0 short & 3.599 & 17.95 & $-5.2\%$ & $+0.123$ & 5.06\% \\
Original, repaired exits & Last price & 3.613 & 18.01 & $-5.2\%$ & $+0.123$ & 5.06\% \\
All rows, original exits (not kept) & $-30\%$ long, 0 short & 3.444 & 16.27 & $-4.7\%$ & $+0.114$ & 6.91\% \\
All rows, repaired exits & $-30\%$ long, 0 short & 3.602 & 17.52 & $-4.7\%$ & $+0.109$ & 5.11\% \\
\textbf{All rows, repaired exits} & \textbf{Last price} & $\mathbf{3.769}$ & $\mathbf{18.19}$ & $\mathbf{-4.6\%}$ & $\mathbf{+0.093}$ & \textbf{5.18\%} \\
\bottomrule
\end{tabular}
\end{table*}

\section{An Audit of the Book}

A review after the first registered run found three problems (Table~\ref{tab:audit}).

\textit{Future information in the universe.} Forecasts were scored only for stocks whose next-month return existed, which is not known at formation and excludes many stocks about to delist. Scoring all rows restores 2,311 name-months (482,241 to 484,552). Trained models are unchanged.

\textit{Exit liquidity.} A delisted name with no current row was filled with the smallest volume among names requested, so exits could cost up to 244\% of notional. The original universe hid most of this; scoring all rows exposes it (IR 3.44), and the repair removes it.

\textit{Missing returns.} A held stock with no return had been booked at $-30\%$ when long and 0 when short. Now both close at the last available price, a return of zero. The book holds 189 such positions: 147 long and 42 short. This is a modelling assumption. Shumway (1997) finds delisting returns are materially negative for performance-related delistings, while others, such as cash mergers, close near the last price, so the truth lies between. \textbf{All of the improvement over the first registered run comes from this choice}: at $-30\%$ the corrected IR is 3.60, below the original 3.63. The $-30\%$ charge is worth 1.99\% a year. Readers should treat 3.60 and 3.77 as a range.

\section{Limitations}

\textit{It is not a clean holdout.} Models are walk-forward, and the last parameter was selected on 2018--2020. But the design passed through 65 trials during which 2021--2026 results for earlier versions had been seen, which gives a deflated Sharpe ratio (Bailey and L\'opez de Prado, 2014) of 0.886, below the usual 0.95. The audit above was also carried out with results in view.

\textit{The shorts are small and hard.} By gross, 91\% of the short book is in micro- and nano-capitalisation stocks (54\% nano), 61\% trades below \$5, and the median short has a \$58m market capitalisation. By the model's dollar-volume tiers, 34\% is hard to borrow or lower; a short is excluded from the lowest tier only on entry, not while held. The model charges 3.4\% a year for borrow and recall, but it has no record of actual loan availability or fee history, and it is here that a real book would most differ.

\textit{Capacity.} Every number here comes from a \$10 million ledger. QUARCC's public paper-fund page presents the same book on a \$1 billion base through a linear $100\times$ rescale, which is for display and is not a capacity estimate. On the development window, a closely related 500-name control already has 34--42\% of short gross beyond one day of 10\% of daily volume and 13--18\% beyond five days, at \$10 million. There, costs grow faster than gross, and net excess return stops rising at about 100--120\% gross. Returns at \$1 billion would be far lower, and we have not estimated by how much.

\textit{Market exposure is not stable.} Full-sample beta is $+0.09$ with a standard error near 0.09, but rolling 12-month beta ranges from $-0.31$ to $+0.99$.

\textit{Other.} All costs are modelled. The live page marks a single 31 August 2026 formation forward and is not a live track record.

\section{Research After the Portfolio}

No later variant replaced the book. Table~\ref{tab:extensions} lists variants from the same research programme, and we report the failures because they show where the edge does and does not lie.

\begin{table*}[t]
\centering
\caption{Variants tested after the portfolio. ``Evaluation'' is January 2021 to August 2026, the same window as Table~\ref{tab:results}; ``Development'' is July 2016 to December 2020. Dashes mean not reported. None replaced the portfolio.}
\label{tab:extensions}
\footnotesize
\setlength{\tabcolsep}{4pt}
\begin{tabular}{llrrrrl}
\toprule
Variant & Window & IR & Vol. & Max DD & Beta & Outcome \\
\midrule
500-name control & Evaluation & 3.478 & 12.4\% & $-4.6\%$ & $+0.098$ & Compliant comparison; excess 47.1\% a year \\
Normal-rank score blend & Evaluation & 3.555 & --- & --- & --- & $-2.12\%$ a year vs.\ control ($t = -1.01$) \\
Volatility target 8\% & Evaluation & 3.697 & 8.7\% & $-3.5\%$ & $+0.065$ & $-11.1\%$ a year vs.\ control ($t = -3.58$) \\
Separated $\mu$, $\Sigma$, cost & Evaluation & 2.29 & 29.8\% & $-28.8\%$ & $+0.77$ & 48.7$\times$ growth, but fails risk and neutrality \\
Unified compiler (best) & Development & 1.04 & 7.3\% & $-6.0\%$ & --- & Below control's 1.36; stopped \\
Fractional IR maximiser & Development & 0.24 & 12.7\% & $-12.4\%$ & --- & $z = -2.26$ vs.\ control; stopped \\
LightGBM mean forecast & Development & 0.62 & 31.6\% & $-57\%$ & $+0.33$ & Stopped \\
Direct parametric policy & Evaluation & --- & 0.44\% & $-1.1\%$ & --- & Near cash; Sharpe 0.29, gross 9--18\% \\
\bottomrule
\end{tabular}
\end{table*}

\textit{Rescaling.} Volatility targeting (Moreira and Muir, 2017) lowers risk but lowers return by more. On the development window, at matched risk, the target arm and the control sit on the same frontier: Sharpe ratio peaks near 8\% volatility for both (1.94 and 1.91), so the arm's higher Sharpe ratio over the full-size control (1.77) reflects running smaller.

\textit{Stronger forecasts do not help on their own.} A normal-rank blend removes the unstable calibration, but it trails the control by a statistically insignificant 2.12\% a year and was not adopted. A separated forecast, risk, and cost system compounds faster ($48.7\times$) and fails its risk limits with a beta of $+0.77$.

\textit{Fuller optimisation is worse.} Unified compilers (Boyd et al., 2024), a fractional IR maximiser (Dinkelbach, 1967), and a direct policy (Brandt, Santa-Clara, and Valkanov, 2009) all fell below the control. The direct policy also shows a general point: policies linear in rank-normalised characteristics spread their bets, while this edge is concentrated.

\textit{A joint risk system fails a placebo test.} We also built a joint risk system with cell-calibrated forecasts and a synthetic economy with planted squeezes and delistings. On a real development window (July 2016 to December 2020), size-quintile specific-risk shrinkage gained $+0.92\%$ a month (95\% CI $[-0.62\%, +2.45\%]$, $p = 0.229$; stepdown-adjusted $p = 0.102$; Romano and Wolf, 2005): the same Sharpe ratio with more risk. In the registered synthetic test over 144 months against 16 shuffled-forecast placebos with the same risk, cost, and constraints, the mean difference was $-0.865\%$ a month (95\% CI $[-2.43\%, +0.70\%]$, $p = 0.283$), the system ranked 16th of 16 (Monte Carlo $p = 0.941$), and an anytime-valid e-process (Ramdas et al., 2023) ended at $\log E = -0.46$. In that world, positive returns came from the machinery around the forecast and not from its ordering. It was never run on 2021--2026 data.

\section*{Availability}

The research code and licensed data are not public. Holdings and marks for the paper book are on QUARCC's public paper-fund page.

\vspace{3pt}
\section*{References}
\scriptsize
\begin{list}{}{\setlength{\leftmargin}{1.3em}\setlength{\itemindent}{-1.3em}\setlength{\itemsep}{1.0pt}\setlength{\parsep}{0pt}}
\item Almgren, Robert, Chee Thum, Emmanuel Hauptmann, and Hong Li, 2005, Direct estimation of equity market impact, \textit{Risk} 18, 58--62.
\item Ang, Andrew, Joseph Chen, and Yuhang Xing, 2006, Downside risk, \textit{Review of Financial Studies} 19, 1191--1239.
\item Bailey, David H., and Marcos L\'opez de Prado, 2014, The deflated Sharpe ratio: Correcting for selection bias, backtest overfitting, and non-normality, \textit{Journal of Portfolio Management} 40, 94--107.
\item Boyd, Stephen, Kasper Johansson, Ronald Kahn, Philipp Schiele, and Thomas Schmelzer, 2024, Markowitz portfolio construction at seventy, \textit{Journal of Portfolio Management} 50, 117--160.
\item Brandt, Michael W., Pedro Santa-Clara, and Rossen Valkanov, 2009, Parametric portfolio policies: Exploiting characteristics in the cross-section of equity returns, \textit{Review of Financial Studies} 22, 3411--3447.
\item Ceria, Sebasti\'an, Anureet Saxena, and Robert A. Stubbs, 2012, Factor alignment problems and quantitative portfolio management, \textit{Journal of Portfolio Management} 38, 29--43.
\item Corwin, Shane A., and Paul Schultz, 2012, A simple way to estimate bid-ask spreads from daily high and low prices, \textit{Journal of Finance} 67, 719--760.
\item D'Avolio, Gene, 2002, The market for borrowing stock, \textit{Journal of Financial Economics} 66, 271--306.
\item Diamond, Steven, and Stephen Boyd, 2016, CVXPY: A Python-embedded modeling language for convex optimization, \textit{Journal of Machine Learning Research} 17, 1--5.
\item Dinkelbach, Werner, 1967, On nonlinear fractional programming, \textit{Management Science} 13, 492--498.
\item Goulart, Paul J., and Yuwen Chen, 2024, Clarabel: An interior-point solver for conic programs with quadratic objectives, arXiv preprint arXiv:2405.12762.
\item Grinold, Richard C., and Ronald N. Kahn, 2000, \textit{Active Portfolio Management}, 2nd ed., McGraw-Hill, New York.
\item Gu, Shihao, Bryan Kelly, and Dacheng Xiu, 2020, Empirical asset pricing via machine learning, \textit{Review of Financial Studies} 33, 2223--2273.
\item Jensen, Theis I., Bryan Kelly, and Lasse Heje Pedersen, 2023, Is there a replication crisis in finance?, \textit{Journal of Finance} 78, 2465--2518.
\item Moreira, Alan, and Tyler Muir, 2017, Volatility-managed portfolios, \textit{Journal of Finance} 72, 1611--1644.
\item Ramdas, Aaditya, Peter Gr\"unwald, Vladimir Vovk, and Glenn Shafer, 2023, Game-theoretic statistics and safe anytime-valid inference, \textit{Statistical Science} 38, 576--601.
\item Romano, Joseph P., and Michael Wolf, 2005, Stepwise multiple testing as formalized data snooping, \textit{Econometrica} 73, 1237--1282.
\item Romano, Yaniv, Evan Patterson, and Emmanuel J. Cand\`es, 2019, Conformalized quantile regression, \textit{Advances in Neural Information Processing Systems} 32, 3543--3553.
\item Shumway, Tyler, 1997, The delisting bias in CRSP data, \textit{Journal of Finance} 52, 327--340.
\item Taleb, Nassim Nicholas, 2020, \textit{Statistical Consequences of Fat Tails}, STEM Academic Press.
\end{list}

\end{document}
